\documentclass{article}
\usepackage[utf8]{inputenc}
\usepackage[T1]{fontenc}
\usepackage[margin=2cm]{geometry}
\usepackage{hyperref}
\sloppy
\begin{document}
\section*{Historical source}
Encyclopaedia Britannica scan, source row 251. Parse 3.5 extracted the text. Super 3.5 VL corrected it from an enlarged crop of each block. It can still contain recognition errors.
Page headers, footers, and page numbers are removed.

\subsection*{Untitled continuation (E0007)}
\textit{Continues an entry from an earlier page.}

Switzerland.

parts of the country that border on Italy or France, a corrupt French or Italian prevails. The two predominant religions are Calvinism and Popery. Of the former are the cantons of Zurich and Berne, the town of St Gall, Geneva, Mühlhausen, and Biel, the principality of Neufchatel, the greater part of Bail, Schaffhausen, the country of the Grisons, the Thurgau, Toggenburg, Glaris, and the Rhine valley; the frontiers of Appenzell, with a small part of Solothurn, and some places in the countries of Baden and Sargans. The rest of the Swiss cantons, allies and dependents, are Popish. For the education of youth there is an university at Basel, and academies at Zurich, Berne, Lausanne, and Geneva, besides gymnasiums and écoles illutres, both in the Popish and Protestant cantons. There are also societies among them for the improvement of the German language and the sciences.

The principal manufactures are snuff and tobacco, linen
of several sorts, lace, thread, silk, and worsted flockings,
neckcloths, cotton stuffs, gloves, handkerchiefs, silks of several
sorts, gold and silver brocades, a variety of woollen manu-
factures, hats, paper, leather of all sorts, earthen wares, por-
celain, toys, watches, clocks, and other hardwares, \&c.
The trade of Switzerland is greatly promoted by many na-
vigable lakes and rivers. In some of the above manufactures,
and in cheese, butter, sheep, horses, black cattle, hides, and
flax, the exports are considerable; and as the imports are
chiefly grain and salt, with some American and Asiatic goods,
there is probably a large balance in their savour. In some
parts of Switzerland dress is restrained by sumptuary laws.

The public revenues are in general very inconsiderable,
arising chiefly from the usual regalia, appropriated every-
where to the foreigne, the demesnes, and public granaries,
voluntary contributions, the sale of salt, and a land-tax; in
the Protestant cantons, from the church-lands also that were
feized at the Reformation. Except in Zurich, Berne, Basel,
and Schaffhausen, where the people are more industrious,
have a greater trade, and are richer than in the others, they
defray the ordinary charges, and that is all.

The cantons never keep any standing troops, except for a few garrisons; but their militia is reckoned to be the best regulated of any in Europe. Every male from 16 to 60 is enrolled, and about one-third of them regimented. They must all provide themselves with arms, clothing, and accoutrements, and appear on the stated days for exercise; and the several cantons and districts must be furnished with a sufficient train of artillery, and all the other implements of war. The Swissers of the several cantons are allowed to engage in the service of such foreign princes and states as are in alliance with those cantons, or with whom they have made a previous agreement. Such states, paying an annual subsidy to the respective cantons, are allowed to make levies. Every man enlists voluntarily, and for what number of years he pleases; at the expiration of which he is at liberty to return home. A great many thus always returning from foreign service, Switzerland is never unprovided with able and experienced officers and soldiers. With respect to their character, they are a brave, honest, hospitable, hardy people; very true to their engagements, friendly, and humane. In short, there is not a people in Europe whose national character is better. In their persons they are generally tall, robust, and well-made; but their complexions are none of the best, and those that live in the neighbourhood of the mountains are subject to wenis. The women are said to be generally handsome and well-shaped, feeble and modest, yet frank, gay, and agreeable in conversation. Few of the peasants are miserably poor; many of them are rich, especially in the Protestant cantons, and that of Berne in particular.

\subsection*{SWIVELS (E0008)}
SWIVELS, a kind of ring made to turn round in a
staple, or other ring. These are used when a ship lies at
her moorings; also in tethers for cattle, that they may turn
round without unwarping the tether.

Swivel-Cannon, is a small piece of artillery belonging to
a ship of war, which carries a shot of half a pound, and is
fixed in a socket on the top of the ship's side, stern, or bow,
and also in her tops. The trunnions of this piece are con-
tained in a fort of iron croutch, of which the lower end ter-
minates in a cylindrical pivot resting in the socket, so as to
support the weight of the cannon. The socket is bored in
a strong piece of oak, reinforced with iron hoops, in order
to enable it to sustain the recoil. By means of this frame,
which is called the swivel, and an iron handle on its cas-
cible, the gun may be directed by the hand to any object.
It is therefore very necessary in the tops, particularly when
loaded with musket-balls, to fire down on the upper decks
of the adversary in action.

\subsection*{SWOONING (E0009)}
SWOONING. See Manchiss, n° 273.

\subsection*{SWORD (E0010)}
SWORD, an offensive weapon worn at the side, and serving either to cut or stab. Its parts are, the handle, guard, and blade; to which may be added the bow, scabbard, pummel, \&c.

Sword of State, which is borne before the king, lords,
and governors of counties, cities, or boroughs, \&c. For
or before the king, it ought to be carried upright; the hilt
as low as the bearer's waist, the blade up between his eyes.
For or before a duke, the blade must decline from the head,
and be carried between the neck and the right shoulder.
For or before an earl, the blade is to be carried between the
point of the shoulder and the elbow: and for or before a ba-
ron, the blade is to be borne in the bend of the arm. This
ceremonial form no less denotes the dignity of a governor
than the coronet set on his coat of arms.

than the coronet set on his coat of arms.

Sword-Fish. See Xiphias.

SWORN, promissory (forensic).

SWORN BROTHERS (frateres jurati), persons who, by mutual oath, covenanted to share each others fortune. Formerly, in any notable expedition to invade and conquer an enemy's country, it was the custom for the more eminent soldiers to engage themselves by reciprocal oaths to share the rewards of their service. This practice gave occasion to the proverb of sworn brothers or brethren in iniquity, because of their dividing plunder and spoil.

SYCAMORE-tree, in botany. See Acer.
SYCOPHANT, an appellation given by the ancient A

\subsection*{SYCOPHANT (E0011)}
SYCOPHANT, an appellation given by the ancient A-
thenians to those who informed of the exportation of figs
contrary to law; and hence it is still used in general for all in-
formers, parasites, flatterers, cheats, \&c.

SYDENHAM (Dr. Thomas), an excellent English physician, was the son of William Sydenham of Winford Eagle in Dorsetshire, and was born there about the year 1624. He studied at Magdalen-hall, Oxford; but left that university when Oxford was garrisoned for king Charles I. and went to London: where, becoming acquainted with Dr Thomas Cox, an eminent physician, that gentleman persuaded him to apply himself to the study of physic; accordingly, after the garrison was delivered up to the parliament, he retired again to Magdalen-hall, entered on the study of medicine, and in 1648 was created bachelor of physic. Soon after, he was made a fellow of All-Souls college, and continued there several years: when, leaving the university, he settled at Westminster, became doctor of his faculty at Cambridge; grew famous for his practice; and was the chief physician in London from the year 1660 to 1670; at which period he began to be disabled by the gout. He died in 1689. His works are highly esteemed both at home and abroad. He was famous for his cool regimen in the small-pox; for giving the bark after the paroxysm in agues; and


\section*{Source record}
\begin{description}
\item[Scan ID] \texttt{britannica:309f765b1ecc83e0a7f7762e44b4a927a7555443:train:251}
\item[Image file] \texttt{page-0251-69952e7c4a0f.jpg}
\item[Image SHA-256] \texttt{69952e7c4a0f40267d357d04308f674155b3ab206ce8f2f4887124ce481f8c1d}
\item[Dataset] NationalLibraryOfScotland/encyclopaedia\_britannica\_illustrated at revision \texttt{309f765b1ecc83e0a7f7762e44b4a927a7555443} (CC0-1.0)
\item[Extraction] model \texttt{nvidia/nemotron-3.5-parse}, raw reply \texttt{extract/raw/page-0251-69952e7c4a0f.json}
\item[Raw reply SHA-256] \texttt{c4807f51f4543e1a7451c1d7a07fcb6aacb326540278d6ea152068cd4bbbf1bb}
\item[Correction] model \texttt{nvidia/Nemotron-3.5-Super-VL}, block crops, raw replies \texttt{correct/raw/crop\_nothink/page-0251-69952e7c4a0f-b*.json}; parse\_fallback: 1, super\_vl: 14
\item[Running header] SWI [ 248] SYD
\end{description}
\end{document}
